\documentclass[10pt]{article}
\usepackage[utf8]{inputenc}
\usepackage[T1]{fontenc}

% --- Configuración de márgenes mejorada para una sola columna ---
\usepackage[a4paper, top=2.5cm, bottom=2.5cm, left=3cm, right=3cm]{geometry}

\usepackage{listings}
\usepackage{xcolor}
\usepackage{titlesec}
\usepackage{hyperref}
\usepackage{fancyhdr}
\usepackage{multicol}
\usepackage{accsupp} 

\hypersetup{
    colorlinks=true,
    linkcolor=black,
    urlcolor=blue
}

% ---------------------------------------------------------------------------
% PARCHES SEGUROS PARA COPIAR Y PEGAR (SIN NÚMEROS Y CON ESPACIOS)
% ---------------------------------------------------------------------------

% 1. Ocultar los números de línea al copiar
\newcommand{\noncopynumber}[1]{%
    \BeginAccSupp{method=plain,ActualText={}}#1\EndAccSupp{}%
}

% 2. Forzar espacios físicos reales al copiar (Método nativo y seguro sin macros rotas)
\makeatletter
\let\orig@lst@outputspace\lst@outputspace
\def\lst@outputspace{%
    \BeginAccSupp{method=hex,unicode,ActualText=0020}%
    \orig@lst@outputspace
    \EndAccSupp{}%
}
\makeatother

% ---------------------------------------------------------------------------
% listings style - Configuración estable para Notebook de CP
% ---------------------------------------------------------------------------
\lstdefinestyle{cppstyle}{%
  language=C++,
  basicstyle=\ttfamily\scriptsize,
  keywordstyle=\color{blue!70!black}\bfseries,
  commentstyle=\color{gray}\itshape,
  stringstyle=\color{green!50!black},
  numberstyle=\tiny\color{gray}\noncopynumber,
  numbers=left,
  stepnumber=1,
  numbersep=4pt,
  frame=single,
  rulecolor=\color{black!25},
  backgroundcolor=\color{gray!5},
  breaklines=true,
  breakatwhitespace=false,
  tabsize=4,
  keepspaces=true,           
  columns=fixed,             
  showstringspaces=false,
  showtabs=false,
  xleftmargin=1.5em,
  framexleftmargin=1.5em
}
\lstset{style=cppstyle}

% Section formatting
\titleformat{\section}{\large\bfseries\scshape}{}{0em}{}[\titlerule]
\titleformat{\subsection}{\normalsize\bfseries}{}{0em}{}

% Header/footer
\pagestyle{fancy}
\fancyhf{}
\lhead{\textsc{Competitive Programming Notebook}}
\rhead{\thepage}
\renewcommand{\headrulewidth}{0.4pt}

\begin{document}

% -------- COVER PAGE --------
\begin{titlepage}
    \centering
    \vspace*{\fill}
    {\Huge\bfseries Competitive Programming Notebook}\\[1em]
    {\large Team: amamemie}\\[0.5em]
    {\normalsize \today}
    \vspace*{\fill}
    \pagenumbering{Gobble} % Apaga la numeración en la portada para evitar el warning de duplicados
\end{titlepage}

% -------- BLANK PAGE --------
\newpage
\thispagestyle{empty}
\mbox{}

% -------- TABLE OF CONTENTS --------
\newpage
\pagenumbering{roman} % Usa números romanos pequeños para el índice (i, ii...)
\tableofcontents
\newpage

\pagenumbering{arabic} % Reinicia a números normales (1, 2, 3...) desde el primer código

\section{BIT MANIPULATION}
\subsection{\texorpdfstring{Binary representation}{Binary representation}}
\lstinputlisting{./Bit_Manipulation/Binary_Representation.cpp}

\section{DATA STRUCTURES}
\subsection{\texorpdfstring{Disjoint set union}{Disjoint set union}}
\lstinputlisting{./Data Structures/Disjoint_Set_Union.cpp}
\subsection{\texorpdfstring{Segment tree}{Segment tree}}
\lstinputlisting{./Data Structures/Segment_Tree.cpp}
\subsection{\texorpdfstring{Sparse table}{Sparse table}}
\lstinputlisting{./Data Structures/Sparse_Table.cpp}

\section{FLOW}
\subsection{\texorpdfstring{Dinic}{Dinic}}
\lstinputlisting{./Flow/Dinic.cpp}

\section{GEOMETRY}
\subsection{\texorpdfstring{Convex hull}{Convex hull}}
\lstinputlisting{./Geometry/Convex_Hull.cpp}
\subsection{\texorpdfstring{Point}{Point}}
\lstinputlisting{./Geometry/Point.cpp}

\section{GRAPH ALGORITHMS}
\subsection{\texorpdfstring{Bfs}{Bfs}}
\lstinputlisting{./Graph Algorithms/Bfs.cpp}
\subsection{\texorpdfstring{Dfs}{Dfs}}
\lstinputlisting{./Graph Algorithms/Dfs.cpp}
\subsection{\texorpdfstring{Dijkstra}{Dijkstra}}
\lstinputlisting{./Graph Algorithms/Dijkstra.cpp}
\subsection{\texorpdfstring{Kruskal mst}{Kruskal mst}}
\lstinputlisting{./Graph Algorithms/Kruskal_MST.cpp}
\subsection{\texorpdfstring{Topological sort kahn algorithm}{Topological sort kahn algorithm}}
\lstinputlisting{./Graph Algorithms/Topological_Sort_Kahn_Algorithm.cpp}

\section{MATH}
\subsection{\texorpdfstring{Binary exponentiation}{Binary exponentiation}}
\lstinputlisting{./Math/Binary_Exponentiation.cpp}
\subsection{\texorpdfstring{Modular inverse}{Modular inverse}}
\lstinputlisting{./Math/Modular_Inverse.cpp}
\subsection{\texorpdfstring{Pollard rho rabin miller}{Pollard rho rabin miller}}
\lstinputlisting{./Math/Pollard_Rho_Rabin_Miller.cpp}
\subsection{\texorpdfstring{Sieve of eratosthenes}{Sieve of eratosthenes}}
\lstinputlisting{./Math/Sieve_of_eratosthenes.cpp}

\section{STRINGS}
\subsection{\texorpdfstring{Hashing}{Hashing}}
\lstinputlisting{./Strings/Hashing.cpp}
\subsection{\texorpdfstring{Kmp}{Kmp}}
\lstinputlisting{./Strings/KMP.cpp}

\section{TREE ALGORITHMS}
\subsection{\texorpdfstring{Binary lifting}{Binary lifting}}
\lstinputlisting{./Tree Algorithms/Binary_Lifting.cpp}
\subsection{\texorpdfstring{Euler tour}{Euler tour}}
\lstinputlisting{./Tree Algorithms/Euler_Tour.cpp}
\subsection{\texorpdfstring{Lowest common ancestor}{Lowest common ancestor}}
\lstinputlisting{./Tree Algorithms/Lowest_Common_Ancestor.cpp}

\end{document}
